% GENERATED FILE. Edit canonical lessons, then run npm run generate:books.
% canonical-source-hash: fnv1a64:2c33bd6daac95627
% canonical-lessons: ZH-C346-kaolu, ZH-C346-guandian, ZH-C346-ganxiang, ZH-C346-yiwei, ZH-C346-guji

\chapter{Thinking It Over}
\label{ch:zh-thinking-it-over}
\hlchaptermodality{346}

\begin{chapteropening}
\textbf{By the end of this chapter:} \emph{I can say I am thinking something over, give a point of view and say I reckon so.}

Complete the last lesson of chapter 346: I can say I am thinking something over, give a point of view and say I reckon so.
\end{chapteropening}

\section[kǎolǜ]{kǎolǜ — to think something over; to consider}

\label{lesson:ZH-C346-kaolu}



\begin{quote}
Before the new one: say the Mandarin for online friend, then the Mandarin for postcode.
\end{quote}

\subsection*{You'll want to know: kǎolǜ}
\emph{kǎolǜ} — \textquotedblleft{}to think something over; to consider\textquotedblright{}.

\begin{quote}
\emph{wǒ zài kǎolǜ qù bu qù Zhōngguó} — I'm thinking over whether to go to China.
\end{quote}

Its characters come later; for now it is a word for the ear.

\begin{cousinweb}
\emph{kǎolǜ} is built from \emph{kǎo} and \emph{lǜ}, both inherited from Middle Chinese, \emph{khawX} and \emph{ljoH} (Guangyun, in Baxter's transcription), which glosses the first 'to check' and 'to complete', and the second 'to think'.
\end{cousinweb}

\begin{grammarlens}[title={What you've built}]
One more word for giving reasons.
\end{grammarlens}

\subsection*{Your turn}
\begin{itemize}
\raggedright
  \item \emph{Say it:} \emph{kǎolǜ}
  \item \emph{Say it:} \emph{kǎolǜ}, once more
  \item \emph{Say it:} \emph{wǒ zài kǎolǜ qù bu qù Zhōngguó} — I'm thinking over whether to go to China.
  \item \emph{Recall:} say the Mandarin for online friend, then the Mandarin for postcode, then say \emph{kǎolǜ} again
\end{itemize}

\begin{tcolorbox}[breakable,colback=teal!4,colframe=teal!35!black,title={Before you move on}]
What does kǎolǜ mean? (\textquotedblleft{}To think something over; to consider\textquotedblright{}.) Say it once more.
\end{tcolorbox}
\section[guāndiǎn]{guāndiǎn — point of view; standpoint}

\label{lesson:ZH-C346-guandian}



\begin{quote}
Before the new one: say the Mandarin for postcode, then the Mandarin for to think something over.
\end{quote}

\subsection*{You'll want to know: guāndiǎn}
\emph{guāndiǎn} — \textquotedblleft{}point of view; standpoint\textquotedblright{}.

\begin{quote}
\emph{wǒ juéde zhège guāndiǎn hěn yǒu yìsi} — I think that's an interesting point of view.
\end{quote}

Its characters come later; for now it is a word for the ear.

\begin{cousinweb}
\emph{guāndiǎn} is built from \emph{guān} and \zh{点} 'a point', both inherited from Middle Chinese, \emph{kwan} and \emph{temX} (Guangyun, in Baxter's transcription), which glosses the first 'to view' and the second 'a dot'.
\end{cousinweb}

\begin{grammarlens}[title={What you've built}]
One more word for giving reasons.
\end{grammarlens}

\subsection*{Your turn}
\begin{itemize}
\raggedright
  \item \emph{Say it:} \emph{guāndiǎn}
  \item \emph{Say it:} \emph{guāndiǎn}, once more
  \item \emph{Say it:} \emph{wǒ juéde zhège guāndiǎn hěn yǒu yìsi} — I think that's an interesting point of view.
  \item \emph{Recall:} say the Mandarin for postcode, then the Mandarin for to think something over, then say \emph{guāndiǎn} again
\end{itemize}

\begin{tcolorbox}[breakable,colback=teal!4,colframe=teal!35!black,title={Before you move on}]
What does guāndiǎn mean? (\textquotedblleft{}Point of view; standpoint\textquotedblright{}.) Say it once more.
\end{tcolorbox}
\section[\texorpdfstring{\zh{感想} (gǎnxiǎng)}{感想 (gǎnxiǎng)}]{\zh{感想} (gǎnxiǎng) — impressions afterwards; thoughts}

\label{lesson:ZH-C346-ganxiang}



\begin{quote}
Before the new one: say the Mandarin for to think something over, then the Mandarin for a point of view.
\end{quote}

\subsection*{You'll want to know: \zh{感想}}
\textbf{\zh{感想}} — \emph{gǎnxiǎng} — \textquotedblleft{}impressions afterwards; thoughts\textquotedblright{}.

\begin{quote}
\textbf{\zh{看了这本书你有什么感想}} \emph{kàn le zhè běn shū nǐ yǒu shénme gǎnxiǎng} — What did you make of the book once you'd read it?
\end{quote}

Two third tones meet in this word: gǎn is said with a rising tone before xiǎng, and the spelling keeps both marks.

\begin{cousinweb}
\textbf{\zh{感想}} (\emph{gǎnxiǎng}) is built from \zh{感} 'to feel' and \zh{想} 'to think', both inherited from Middle Chinese, \emph{komX} and \emph{sjangX} (Guangyun, in Baxter's transcription), which glosses the first 'to move' and the second 'to think'.
\end{cousinweb}

\begin{grammarlens}[title={What you've built}]
One more word for giving reasons.
\end{grammarlens}

\subsection*{Your turn}
\begin{itemize}
\raggedright
  \item \emph{Say it:} \emph{gǎnxiǎng}
  \item \emph{Say it:} \emph{gǎnxiǎng}, once more
  \item \emph{Say it:} \emph{kàn le zhè běn shū nǐ yǒu shénme gǎnxiǎng} — What did you make of the book once you'd read it?
  \item \emph{Recall:} say the Mandarin for to think something over, then the Mandarin for a point of view, then say \emph{gǎnxiǎng} again
\end{itemize}

\begin{tcolorbox}[breakable,colback=teal!4,colframe=teal!35!black,title={Before you move on}]
What does \zh{感想} mean? (\textquotedblleft{}Impressions afterwards; thoughts\textquotedblright{}.) Say it once more.
\end{tcolorbox}
\section[\texorpdfstring{\zh{以为} (yǐwéi)}{以为 (yǐwéi)}]{\zh{以为} (yǐwéi) — to wrongly assume; to take it that (often wrongly)}

\label{lesson:ZH-C346-yiwei}



\begin{quote}
Before the new one: say the Mandarin for a point of view, then the Mandarin for your impressions afterwards.
\end{quote}

\subsection*{You'll want to know: \zh{以为}}
\textbf{\zh{以为}} — \emph{yǐwéi} — \textquotedblleft{}to wrongly assume; to take it that (often wrongly)\textquotedblright{}.

\begin{quote}
\textbf{\zh{我以为你不来了}} \emph{wǒ yǐwéi nǐ bù lái le} — I thought you weren't coming.
\end{quote}

\begin{cousinweb}
\textbf{\zh{以为}} (\emph{yǐwéi}) is built from \zh{以} 'to take' and \zh{为} 'to make', both inherited from Middle Chinese, \emph{yiX} and \emph{hjwe} (Guangyun, in Baxter's transcription), which glosses the first 'to use' and lists the second as the everyday spelling of a word glossed 'to make'. How the pair came to mean a belief that turns out wrong is not settled.
\end{cousinweb}

\begin{grammarlens}[title={What you've built}]
One more word for giving reasons.
\end{grammarlens}

\subsection*{Your turn}
\begin{itemize}
\raggedright
  \item \emph{Say it:} \emph{yǐwéi}
  \item \emph{Say it:} \emph{yǐwéi}, once more
  \item \emph{Say it:} \emph{wǒ yǐwéi nǐ bù lái le} — I thought you weren't coming.
  \item \emph{Recall:} say the Mandarin for a point of view, then the Mandarin for your impressions afterwards, then say \emph{yǐwéi} again
\end{itemize}

\begin{tcolorbox}[breakable,colback=teal!4,colframe=teal!35!black,title={Before you move on}]
What does \zh{以为} mean? (\textquotedblleft{}To wrongly assume; to take it that (often wrongly)\textquotedblright{}.) Say it once more.
\end{tcolorbox}
\section[gūjì]{gūjì — to reckon; to make a rough guess}

\label{lesson:ZH-C346-guji}



\begin{quote}
Before the new one: say the Mandarin for your impressions afterwards, then the Mandarin for to wrongly assume.
\end{quote}

\subsection*{You'll want to know: gūjì}
\emph{gūjì} — \textquotedblleft{}to reckon; to make a rough guess\textquotedblright{}.

\begin{quote}
\emph{yīnwèi xiàyǔ suǒyǐ wǒ gūjì tā bú huì lái} — It's raining, so I reckon he won't come.
\end{quote}

Its characters come later; for now it is a word for the ear.

\begin{cousinweb}
\emph{gūjì} is built from \emph{gū} and \emph{jì}, both inherited from Middle Chinese, \emph{kuX} and \emph{kejH} (Guangyun, in Baxter's transcription), which glosses the first 'a market tax' and the second 'to reckon up'. How the first came to mean estimating is not settled.
\end{cousinweb}

\begin{grammarlens}[title={What you've built}]
One more word for giving reasons.
\end{grammarlens}

\subsection*{Your turn}
\begin{itemize}
\raggedright
  \item \emph{Say it:} \emph{gūjì}
  \item \emph{Say it:} \emph{gūjì}, once more
  \item \emph{Say it:} \emph{yīnwèi xiàyǔ suǒyǐ wǒ gūjì tā bú huì lái} — It's raining, so I reckon he won't come.
  \item \emph{Recall:} say the Mandarin for your impressions afterwards, then the Mandarin for to wrongly assume, then say \emph{gūjì} again
\end{itemize}

\begin{tcolorbox}[breakable,colback=teal!4,colframe=teal!35!black,title={Before you move on}]
What does gūjì mean? (\textquotedblleft{}To reckon; to make a rough guess\textquotedblright{}.) Say it once more.
\end{tcolorbox}
